%% --------------------------------------------------------------------------
\section{Query Implementation Methodology}
\label{sec:eval:query-impl}

\project{} exposes a C++ cursor API but provides no query engine or declarative query language.
Neo4j, ArangoDB, and NeuG accept LDBC SNB queries written in Cypher or AQL; their query engines parse the declarative text, optimize join order and access paths, and execute the resulting plan.
Evaluating \project{} against these systems requires hand-written C++ implementations of each benchmark query: imperative code that performs the same traversals, filters, and aggregations that a query engine would derive from the declarative specification.

We used Claude (Anthropic) to generate these implementations in a three-stage pipeline: schema interface generation, query logic synthesis, and cross-system verification.
The LLM received the LDBC SNB Interactive v1 query specifications and the \project{} source code (header files defining the \code{GraphBase} API, cursor classes, and property schema namespaces) as context.

\paragraph{Schema interface generation.}
The first stage produced the mapping between the LDBC SNB data model and \project{}'s storage primitives.
This mapping has three parts: vertex type encoding, property schemas, and edge-type routing.

LDBC SNB defines 11 concrete vertex types (Person, Post, Comment, Forum, Tag, TagClass, City, Country, Continent, Company, University) and 20 edge types.
\project{} encodes vertex type in the top 8~bits of a 64-bit vertex ID (\autoref{arch:label-encoding}); the LLM assigned each LDBC type a \code{VertexType} enum value (\code{VT\_PERSON}=0 through \code{VT\_UNIVERSITY}=10) and generated the bulk loader that remaps LDBC's arbitrary large integer IDs to compact typed IDs via per-type hash maps.
The Place and Organisation hierarchies (City/Country/Continent and Company/University) each receive distinct enum values rather than a shared one, because edges such as \textsc{studyAt} (Person$\to$University) and \textsc{workAt} (Person$\to$Company) would otherwise produce ambiguous \code{(src\_type, dst\_type)} pairs.

For each vertex type, the LLM generated a property schema namespace, a C++ struct defining the binary layout of the property blob with fixed byte offsets and typed accessor functions (e.g., \code{SNBPersonSchema}: 153~bytes covering \code{creationDate}, \code{birthday}, \code{gender}, \code{firstName}, \code{lastName}, \code{browserUsed}, \code{locationIP}, \code{country\_id}, and \code{city\_id}).
It also designed the \wt{} column-group assignments for columnar mode: each vertex type's properties are partitioned into column groups that separate fixed-size, scan-friendly columns from variable-length ones.
For example, Person properties split into four groups: \emph{temporal} (\code{creationDate}, \code{birthday}; 16~bytes per row), \emph{name} (\code{firstName}, \code{lastName}, \code{gender}), \emph{contact} (\code{browserUsed}, \code{locationIP}), and \emph{location} (\code{country\_id}, \code{city\_id}).
Post and Comment properties split into \emph{temporal} (\code{creationDate}, \code{length}; 12~bytes) and \emph{content} (\code{tag}, \code{content}; variable-length).
Multi-valued attributes (Person email addresses and spoken languages) are stored in secondary \wt{} tables with composite keys \code{(person\_id, index)}.

Edge property schemas follow a similar pattern.
Five of the 20 edge types carry properties: \textsc{knows}, \textsc{likes}, and \textsc{hasMember} each store a \code{creationDate} (8~bytes); \textsc{studyAt} stores \code{classYear} (4~bytes); \textsc{workAt} stores \code{workFrom} (4~bytes).
The remaining 15 edge types are structural-only and use a 1-byte sentinel value.
Edge types are disambiguated by endpoint vertex types: \textsc{knows} connects \code{VT\_PERSON}$\to$\code{VT\_PERSON}, \textsc{likes} connects \code{VT\_PERSON}$\to$\code{VT\_POST} or \code{VT\_COMMENT}, and so on.
The LLM generated the routing functions: \code{get\_edge\_properties} calls access the appropriate property tables by first extracting the \code{(src\_type, dst\_type)} pair from the given vertex IDs.

\paragraph{Query logic synthesis.}
The second stage translated each LDBC query specification into a C++ function operating on \project{}'s cursor API.
For each query, the LLM received the specification's input parameters, output schema, sort order, and limit, along with the schema interfaces from the first stage.
It produced imperative traversal code: topology cursors (\code{get\_out\_nodes\_id}, \code{get\_in\_nodes\_id}) for structural hops, property cursors (\code{NodePropCursor}, \code{EdgePropCursor}) for attribute access, and standard-library containers for intermediate results.
Point queries (R1, X1) reduce to a single \code{get\_node\_properties} call and schema deserialization.
Multi-hop queries (IC-3, IC-9) expand to nested loops over neighbor lists with property-predicate filters.
BI queries (BI-1, BI-12) use sequential column-group cursor scans over all vertices of a type.
The generated code totals approximately 2{,}500 lines of C++ across 18~query functions.

\paragraph{Verification and conformance.}
% Each query we implemented in \project{} was validated against Neo4j by running both systems on the same LDBC SNB SF-3 dataset and comparing results field-by-field.
We validated each generated query by comparing the output of \project{} to that produced by Neo4j.
A Python script (\code{validate\_results.py}) invokes \project{}'s benchmark binary in validation mode and the corresponding Cypher query on Neo4j, parses both outputs into sorted JSON records, and reports any discrepancies in values, sort order, or result-set cardinality.
Validation runs across all four \project{} storage configurations (Adjacency~List and EdgeKey, each with embedded and columnar properties) to confirm that query semantics are independent of the underlying storage layout.
All 18 query functions must produce identical results across all four configurations before any benchmark timing is collected.

The verification process also served as a conformance check against the LDBC specification.
When the LLM's initial implementation of a query diverged from the spec, the cross-system comparison exposed the error: for instance, the first IC-3 implementation counted common friends between a person and each friend-of-friend candidate, when the specification requires counting each candidate's messages located in two given countries during a date window.
Similarly, the initial IC-5 counted friends per forum rather than posts per forum by friends-of-friends who joined after a cutoff date.
In both cases, the field-by-field comparison against Neo4j's Cypher output flagged the semantic mismatch, and the LLM regenerated the query logic from the specification.
The NeuG Cypher queries were updated in lockstep to maintain identical semantics across all three systems.

The final conformant set covers 14 queries from the LDBC SNB Interactive v1 specification: 4 short reads (SR-1, SR-3, SR-5, SR-7), 7 complex reads (IC-2, IC-3, IC-5, IC-7, IC-8, IC-9, IC-13), and 3 insert operations (INS-1, INS-6, INS-8).
The remaining 4 query functions (R1, X1, A2, A3, X5) are targeted microbenchmarks that isolate individual access patterns (single-vertex property fetches and edge-count aggregates with temporal predicates) that the LDBC complex queries bundle with multi-hop traversal, making it difficult to attribute latency to the property layer alone.

\paragraph{Rationale.}
For a novel storage system without a query engine, the alternative to LLM-assisted generation is manual implementation of each benchmark query, the approach used by prior cursor-based systems such as GraphflowDB~\cite{graphflowdb2021Gupta} and LiveGraph~\cite{zhu2019livegraph}, which publish hand-written query harnesses for their evaluations.
The LLM-assisted approach derives query logic from the same LDBC specification that other systems implement through their query engines.
The conformance errors caught during verification exposed a secondary value of this pipeline: the same cross-system comparison that validates correctness also enforces spec compliance, catching semantic drift that would otherwise go unnoticed in a hand-written harness until benchmark results are questioned.
Cross-system verification against Neo4j ensures that any performance differences reflect storage-layer behavior, not divergent query semantics.
